% 封面页：参赛项目说明书封面
% ============================================

\newcommand{\makecover}{
    \thispagestyle{empty}

    \begin{center}
        % ---------- 标题区（三行居中） ----------
        {\fontsize{22pt}{28pt}\selectfont\heiti
        第十八届（\competitionYear{}年）}

        \vspace{0.6cm}

        {\fontsize{26pt}{32pt}\selectfont\heiti
        北京市大学生物理实验竞赛}

        \vspace{0.6cm}

        {\fontsize{24pt}{30pt}\selectfont\heiti
        参赛项目说明书}

        \vspace{2cm}

        % ---------- 参赛题目与类别选择表 ----------
        {\fontsize{14pt}{18pt}\selectfont\heiti 参赛题目与类别：（请在□处画√）}

        \vspace{1cm}

        % 使用 tabular 构建选择表（与原始格式一致）
        \renewcommand{\arraystretch}{1.6}
        \begin{tabular}{@{}p{12cm}@{}}
            % --- 命题类 ---
            {\heiti 命题类} \\
            \makecheckbox{\topicOne} 题目1：微小位移测量 \\
            \makecheckbox{\topicTwo} 题目2：探究电磁感应现象中的能量转换 \\
            \makecheckbox{\topicThree} 题目3：弱压力测量 \\
            \makecheckbox{\topicFour} 题目4：晶体双折射 \\[6pt]

            % --- 制作类 ---
            {\heiti 制作类} \\
            \makecheckbox{\topicFive} 题目5：实验仪器的制作、改进
                （\makecheckbox{\topicFiveSubA}定量类 \; \makecheckbox{\topicFiveSubB}演示类） \\
            \makecheckbox{\topicSix} 题目6：前沿物理
                （\makecheckbox{\topicSixSubA}前沿文献类 \; \makecheckbox{\topicSixSubB}自研文献类） \\
            \makecheckbox{\topicSeven} 题目7：AI+物理实验 \\[6pt]

            % --- 信息资源类 ---
            {\heiti 信息资源类} \\
            \makecheckbox{\topicEight} 题目8：大学物理教学微视频 \\
            \makecheckbox{\topicNine} 题目9：物理教学资源开发
                （\makecheckbox{\topicNineSubA}媒体资源类 \; \makecheckbox{\topicNineSubB}仿真程序类） \\
        \end{tabular}

        \vspace{2.5cm}

        % ---------- 组委会信息 ----------
        {\fontsize{13pt}{16pt}\selectfont\songti 北京市大学生物理实验竞赛组委会制}

        \vspace{0.5cm}

        {\fontsize{13pt}{16pt}\selectfont\songti \competitionYear{}年6月}

    \end{center}

    \clearpage
}
